% ============================================================================
% Chapter 10 --- Storage and data classification
% ----------------------------------------------------------------------------
% Purpose : object/block storage, classification-aware access, BYOK.
%           Belongs to the Execution plane (D5.1).
% Page target: ~4 pages.
% Dependencies: Ch. 5 (S3 API, KMIP, NFS v4.1), Ch. 7 (policy decisions
%               for classification-driven access), Ch. 9 (Ceph-CSI as
%               reference --- cross-link from D9.2), Ch. 13 (cryptographic
%               services for KMS).
% Mirrors layer chapter template established in Ch. 6 (D6.5).
% ----------------------------------------------------------------------------
% Decisions registered in _coherence-EN.md:
%   D10.1 minimum compliance set for the storage layer
%   D10.2 reference instantiation = Ceph + Vault Transit + label-based access
%   D10.3 alternatives = NetApp/KMIP, hybrid on-prem + cloud with strict BYOK
%   D10.4 storage classification scheme is institutional, not prescribed
%   D10.5 storage sovereignty scores
% ============================================================================

\chapter{Storage and data classification}
\label{ch:storage}

\begin{caveatbox}[Dependencies]
\noindent This chapter applies the layer template of
\trchap{ch:identity}{} to the storage layer. It consumes
classification-driven policy decisions from \trchap{ch:policy}{}, it
binds to compute workloads through the CSI plugins referenced in
\trchap{ch:compute}{}, and it relies on the cryptographic services
of \trchap{ch:crypto}{}. The layer belongs to the \emph{Execution
plane}.
\end{caveatbox}

The storage layer is where the data dimension of sovereignty is
most directly tested. A network configuration can be revised in
weeks; a misclassified or mislocated data corpus is rarely revised at
all. The architectural commitment of this layer is consequently
modest in scope but strict in observance: standard interfaces,
explicit classification, institutional key custody, and an
exit-data export procedure that is exercised, not assumed.

%-----------------------------------------------------------------------------
\section{Functional role and interface contract}
\label{sec:st-contract}

The storage layer is responsible for: providing object and block
storage to workloads through standard interfaces; carrying a data
classification label that the policy layer consults; encrypting
data at rest with keys whose custody is governed; permitting the
extraction of institutional data in documented open formats; and
producing access events ingestible by the observability layer.

\begin{contractbox}[Storage layer]
\noindent A compliant storage layer, in the sense of this
architecture, must:

\begin{itemize}[leftmargin=1.5em,nosep]
  \item expose object storage through the S3 API \emph{or} an
    interface broadly compatible with it;
  \item expose shared-file access through NFS~v4.1~\cite{rfc-8881-nfs41}
    or an equivalent standard, where shared-file access is required;
  \item carry, for each stored object or volume, a
    classification label from a documented institutional scheme,
    accessible to \trchap{ch:policy}{};
  \item encrypt data at rest, with key custody held in a system that
    speaks KMIP~\cite{kmip-oasis} \emph{or} an equivalent open
    key-management interface, and that admits a Bring-Your-Own-Key
    arrangement;
  \item produce access events (object created, read, written,
    deleted, exported) ingestible by the observability layer in a
    schema consistent with \trchap{ch:observability}{};
  \item document an exit-data export procedure that delivers stored
    data in a documented open format, and exercise it periodically
    as part of \trchap{ch:conformance}{}.
\end{itemize}
\end{contractbox}

\paragraph{Classification scheme.} The contract requires a documented
classification scheme but does not prescribe one. In practice,
institutions tend to adopt either a Traffic Light Protocol variant,
a domain-specific scheme (for example, those used in medical or
defence-adjacent research), or an institution-specific scheme that
maps to the regulatory regimes the institution operates under. The
choice belongs to the institution; the contract requires that the
choice be documented and that every label used be enumerable.

%-----------------------------------------------------------------------------
\section{Reference instantiation: Ceph + Vault Transit + label-based access}
\label{sec:st-reference}

The reference instantiation combines, as representative open-source
implementations, a distributed storage system such as
Ceph~\cite{ceph} providing both object (RGW) and block (RBD)
interfaces; an envelope-encryption secrets engine with
institutionally-held keys such as HashiCorp Vault~\cite{vault}'s
Transit engine; and label-based access
control implemented through a policy decision integrating the
classification label into the authorisation logic.

\begin{instantiationbox}[Storage layer]
\noindent\textbf{Topology.} Ceph clusters provide block storage to
\trchap{ch:compute}{} workloads via Ceph-CSI (consistent with
\S\ref{sec:co-reference}), object storage via the RGW gateway
exposing the S3 API, and file storage via CephFS where shared-file
semantics are required. Encryption at rest is enabled at the OSD
level, with keys wrapped through Vault Transit; the institutional
key-management policy resolves to either a Vault-internal master
key or, for higher-classification data, an HSM-backed master key
exposed through PKCS\#11. Each object or volume carries a
classification label as metadata; the policy decision for read or
modify access combines the subject attributes from
\trchap{ch:identity}{} with the label, and the decision is logged
through the observability layer.
\end{instantiationbox}

\paragraph{Operational considerations.} Operating Ceph at
institutional scale would typically require dedicated expertise:
the cluster's health and recovery characteristics demand specific
operational discipline that more conventional storage products do
not require. The staffing envelope follows the institution's
established operational staffing profile for distributed storage,
typically a small dedicated team with substantial Ceph familiarity,
often shared with the network and compute teams. The labour market
is broad for Ceph itself and narrower for the integration of Ceph
with classification-aware access control.

\paragraph{Cryptographic posture.} Vault Transit provides envelope
encryption that allows the encryption keys to be rotated without
re-encrypting the data corpus. For data classifications whose
disclosure regime warrants stronger custody, the institution would
substitute a HYOK arrangement, consistent with the institution's own
data-custody sovereignty profile, in which Vault never holds the
cleartext key and cryptographic operations involve the institution's
HSM in the live path. This shifts the data-dimension score upward
at the cost of additional operational complexity, addressed in
\trchap{ch:crypto}{}.

\begin{realworldbox}[Research storage precedents]
\noindent CERN's EOS distributed file system, designed for high-
energy physics datasets, is one of the most extensively documented
research-storage architectures and reflects the operational scale
that physical-sciences research can require. European research
institutions participate in the EUDAT Collaborative Data
Infrastructure (EUDAT~CDI), whose B2 services suite (B2SAFE for
replication, B2SHARE for publication, B2STAGE for compute access,
B2DROP for sync, B2FIND for discovery, B2ACCESS for federated
authorisation, B2HANDLE for persistent identifiers) operates
across a network of more than twenty research and computing
centres in over a dozen European nations. These European research-computing references should
be consulted for their design lessons rather than as direct
templates --- adopters in other regions have analogous national or
regional research-data infrastructures --- since the relative scale
of teaching, administrative and research storage at a given
institution may produce a substantially different proportion than
at a large-scale research-computing environment.
\end{realworldbox}

%-----------------------------------------------------------------------------
\section{Alternative~1 --- Enterprise storage platform with KMIP-compliant HSM}
\label{sec:st-alt1}

\begin{alternativebox}[Enterprise SAN/NAS]
\noindent The institution operates an enterprise storage platform
from one of the established vendors in this market that exposes the
S3 API for object
storage and NFS for shared file access, with encryption at rest
keyed through a KMIP-compliant HSM under institutional control.
Classification metadata is carried at the application level rather
than natively in the storage system.
\end{alternativebox}

\paragraph{Where it adds value.} A mature platform with predictable
performance characteristics, vendor support, and integration with
existing operational tooling. For institutions whose storage is
predominantly transactional rather than research-scale, the
operational simplicity may be substantial.

\paragraph{Where it constrains the institution.} Two recurring
constraints. First, classification metadata typically lives outside
the storage system and must be reconciled with it through
application or integration logic; the reconciliation is a
non-trivial engineering effort that the reference instantiation
avoids. Second, exit-data export from proprietary storage platforms
is variable in cost and contractually framed; the
exit-cost-proportionality indicator of
\trchap{ch:sovereignty-grid}{} should be evaluated explicitly at
procurement.

\paragraph{When it could be the right choice.} An institution with
limited research-storage requirements, a mature operational
relationship with an enterprise storage supplier, and a
sovereignty posture that accepts the trade-offs in exchange for
the operational simplification.

%-----------------------------------------------------------------------------
\section{Alternative~2 --- Hybrid (on-premises + S3-compatible cloud with strict BYOK)}
\label{sec:st-alt2}

\begin{alternativebox}[Hybrid storage with BYOK]
\noindent The institution operates a storage tier under an
institution-controlled deployment custody model (Ceph or enterprise,
self-hosted or otherwise kept under direct institutional control) for
data whose classification requires institutional custody, and
federates additional capacity to an S3-compatible cloud
provider for less-sensitive classes of data, with all cloud-stored
data encrypted under institutionally-held keys (BYOK, with HYOK for
the higher-sensitivity classes that the hybrid posture would
admit).
\end{alternativebox}

\paragraph{Where it adds value.} Capacity flexibility for
research-data corpora whose growth profile would over-provision an
on-premises footprint, and geographic redundancy for data whose
classification permits multi-jurisdiction storage.

\paragraph{Where it constrains the institution.} The data dimension
of \trchap{ch:sovereignty-grid}{} is permanently bounded by the
weakest disclosure regime applicable to any portion of the corpus.
The institutional governance of which data classes are eligible for
the cloud tier becomes a structural concern, addressed through the
policy layer (\trchap{ch:policy}{}) and the data-classification
scheme.

\paragraph{When it could be the right choice.} An institution
whose research storage profile combines a sensitive on-premises core
with a substantial volume of less-sensitive data (publication-stage
artefacts, public datasets), and whose classification governance is
mature enough to enforce the boundary reliably.

%-----------------------------------------------------------------------------
\section{Sovereignty grid applied}
\label{sec:st-sovereignty}

\begin{table}[H]
\centering\small
\caption{Per-dimension sovereignty profile of the storage
instantiations.}
\label{tab:st-sovereignty}
\begin{tabularx}{\textwidth}{%
  >{\hsize=0.6\hsize\linewidth=\hsize\bfseries\raggedright\arraybackslash}X%
  >{\hsize=1.2\hsize\linewidth=\hsize\centering\arraybackslash}X%
  >{\hsize=1.2\hsize\linewidth=\hsize\centering\arraybackslash}X%
  >{\hsize=1.2\hsize\linewidth=\hsize\centering\arraybackslash}X%
  >{\hsize=0.8\hsize\linewidth=\hsize\centering\arraybackslash}X%
}
\toprule
Dimension & Reference (Ceph + Vault) & Enterprise platform + KMIP HSM & Hybrid (on-prem + cloud) \\
\midrule
Data          & 3 Robust      & 3 Robust      & 1--2 Partial  \\
Technical     & 3 Robust      & 1 Partial     & 2 Established \\
Financial     & 2 Established & 1 Partial     & 2 Established \\
Operational   & 2 Established & 3 Robust      & 2 Established \\
Institutional & 3 Robust      & 2 Established & 2 Established \\
\bottomrule
\end{tabularx}
\end{table}

\paragraph{Driving indicators by dimension.} Data: storage location,
key custody, disclosure regime --- this is the dimension on which
the storage layer is most often constrained. Technical: source
visibility and the reconciliation of classification with the
storage system. Financial: cost predictability and exit-cost
proportionality. Operational: depth of the skills market for the
chosen technology. Institutional: control over the storage
roadmap.

\noindent The hybrid pattern's data-dimension score is bounded by
the cloud-tier disclosure regime, even when the on-premises tier
scores Robust. This is the structural reason why classification
governance must be in place before the hybrid pattern is adopted.

%-----------------------------------------------------------------------------
\section{Regulatory anchoring}
\label{sec:st-regulatory}

\noindent The storage layer is the principal anchor of the
data-dimension obligations under the applicable data-protection
regime: security of processing (e.g.\ \gls{gdpr}~Art.~32), data
minimisation through tiering (e.g.\ Art.~5(1)(c)), and the regime's
constraints on international transfers (e.g.\ Art.~44--49), which bind
when storage placement crosses jurisdictions. It is the operational
counterpart of \loc{cross-border-transfer-regime}:
storage-placement constraints, key-custody and per-tenant key
separation are the technical supplementary measures that an
institution can invoke alongside whatever standard contractual
transfer instruments that regime recognises. The layer
operationalises NIS2~Art.~21(2)(c) (backup) and
contributes to 27001~A.8.10 (information deletion), A.8.12
(data leakage prevention) and A.8.13 (information backup). It
contributes to NIST CSF \textsf{PROTECT} (PR.DS --- data
security). Detailed mapping in \trchap{ch:regulatory-mapping}{}
\S\ref{sec:rm-gdpr}, \S\ref{sec:rm-schrems2},
\S\ref{sec:rm-nis2}, \S\ref{sec:rm-iso}.

%-----------------------------------------------------------------------------
\section{Common pitfalls}
\label{sec:st-pitfalls}

\begin{pitfallbox}[Classification drift]
\noindent The classification labels attached to data corpora at
creation time tend to remain unchanged as the data evolves --- as
new fields are added, as additional populations are sampled, as
linkages with other corpora are introduced. The classification
becomes increasingly disconnected from the actual sensitivity. The
mitigation is to schedule a classification review at data-corpus
events (substantial expansion, linkage with another corpus,
publication) and to record the classification history alongside the
data.
\end{pitfallbox}

\begin{pitfallbox}[Key-management lock-in]
\noindent A storage system whose encryption keys can only be managed
through the supplier's own key-management product creates a
structural exit cost: the keys cannot be transferred to another
storage system without re-encrypting the entire corpus. The
mitigation is to require KMIP (or an equivalent open interface) at
procurement, and to verify that the institution's keys can in fact
be operated through a non-supplier key-management product.
\end{pitfallbox}

\begin{pitfallbox}[Opaque snapshots and replicas]
\noindent Storage systems take snapshots and create replicas as part
of normal operation; the location, retention, and access regime of
those snapshots is not always documented to the institution. A
data-dimension audit that ignores snapshots can score a corpus as
Robust on storage location while replicas of it sit elsewhere. The
mitigation is to enumerate snapshots and replicas as part of the
data-dimension self-assessment, and to bring their location under
the same governance as the primary storage.
\end{pitfallbox}

\begin{pitfallbox}[Cross-region replication outside jurisdiction]
\noindent Cloud-tier storage often defaults to multi-region
replication for durability; the regions selected by the supplier
may sit outside the institution's intended jurisdiction. The
mitigation is to declare region restrictions explicitly at
provisioning, to verify them through the access events of
\trchap{ch:observability}{}, and to test the assumption
periodically through deliberate write-and-read patterns that would
surface unexpected replication.
\end{pitfallbox}
